%==============================================================================
% CHAPTER 27: DNA Matching as Wave Interference
%==============================================================================

\chapter{DNA Matching as Wave Interference}
\label{chap:wave_interference}

\epigraph{%
  Young's fringes are not a curiosity. They are the signature of a wave that
  knows where it has been.%
}{Max Born, \textit{Atomic Physics} (1935)}

\begin{chapterbox}
If sequences are oscillations (Chapter~\ref{chap:bounded_phase_space},
Chapters~\ref{chap:shader_spectrometer}) then sequence matching is wave
interference. This chapter derives the Matched Filter Optimality Theorem: the
matched filter cross-correlation is the minimum mean-square error detector for
a query waveform embedded in a noisy target, and is therefore the theoretically
correct method for DNA sequence alignment. Phase-preserving matched filter
analysis achieves 2\,bp peak width; magnitude-only analysis (spectral embedding)
broadens the peak to 100--130\,bp, with the discrepancy explained exactly by
the Nyquist sampling theorem applied to the cross-correlation function.
Validated results from \citet{Sachikonye2025_Interference}: 60/60 perfect
position recovery at 50\% mutation rate; AUC $= 1.000$ at all tested mutation
rates (0\%--50\%); F1 $= 1.000$ at $z$-score $\geq 6$; background statistics
normally distributed under the null, validating $z$-score calibration.
Preprocessing 1\,Mb target: $\approx 530$\,ms; per-query scan 1\,Mb:
$\approx 836$\,ms.
\end{chapterbox}

%------------------------------------------------------------------------------
\section{The Wave Identity for Sequence Matching}
\label{sec:wave_identity}
%------------------------------------------------------------------------------

Chapters~\ref{chap:bounded_phase_space}--\ref{chap:shader_spectrometer}
established the chain: nucleotide sequences are oscillations in cardinal
coordinate space; oscillations are partitions; partitions are categorical
distinctions. The cardinal map $\Cardinal$ (equation~\eqref{eq:cardinal_map})
is a functor from nucleotide sequences to waveforms in $\mathbb{Z}^2$.

\begin{proposition}[Sequence Matching is Wave Interference]
\label{prop:matching_is_interference}
Under the cardinal embedding $\Cardinal$, the problem of locating a query
sequence $q$ within a target sequence $t$ is identical to the problem of
detecting a known waveform $\Cardinal(q)$ embedded in a noisy background
waveform $\Cardinal(t)$.
\end{proposition}

\begin{proof}
The cardinal map $\Cardinal$ is an isometric embedding (distances in sequence
space are preserved by categorical distance $\dC$; see
Definition~\ref{def:categorical_distance}). Under $\Cardinal$, sequence
similarity becomes waveform correlation: two similar sequences produce similar
waveforms in $\mathbb{Z}^2$ (high cross-correlation), while dissimilar sequences
produce orthogonal waveforms (low cross-correlation). The problem of finding
$q$ at offset $k^*$ within $t$---sequence alignment---maps exactly to the
signal detection problem: determine the lag $k^*$ at which $\Cardinal(q)$
maximally correlates with a window of $\Cardinal(t)$. The background of
$\Cardinal(t)$ away from the alignment offset constitutes the noise in the
detection problem. \qed
\end{proof}

The optimal detector for a known signal in noise is the matched filter,
a classical result from signal detection theory. The following section derives
this result in the partition-theoretic framework.

%------------------------------------------------------------------------------
\section{The Matched Filter Cross-Correlation}
\label{sec:matched_filter}
%------------------------------------------------------------------------------

\begin{definition}[Cardinal Cross-Correlation]
\label{def:cross_correlation}
For cardinal-embedded sequences $\mathbf{q} = \Cardinal(q)$ and
$\mathbf{t} = \Cardinal(t)$ of lengths $L$ and $T$ respectively ($T \geq L$),
the \emph{matched filter cross-correlation} is:
\begin{equation}
  \mathrm{xcorr}(q, t)[k] = \mathrm{IFFT}\bigl(\overline{\mathrm{FFT}(\mathbf{q})} \cdot \mathrm{FFT}(\mathbf{t})\bigr)[k],
  \quad k = 0, 1, \ldots, T-1,
  \label{eq:xcorr}
\end{equation}
where $\overline{\,\cdot\,}$ denotes complex conjugate and $\cdot$ denotes
element-wise multiplication.
\end{definition}

\begin{remark}[Two-Channel Extension]
For DNA sequences with two-channel cardinal embedding, the cross-correlation
is computed independently on each channel:
\begin{align}
  \mathrm{xcorr}_x[k] &= \mathrm{IFFT}\bigl(\overline{\mathrm{FFT}(\mathbf{x}_q)} \cdot \mathrm{FFT}(\mathbf{x}_t)\bigr)[k], \\
  \mathrm{xcorr}_y[k] &= \mathrm{IFFT}\bigl(\overline{\mathrm{FFT}(\mathbf{y}_q)} \cdot \mathrm{FFT}(\mathbf{y}_t)\bigr)[k],
\end{align}
and combined by taking $\mathrm{xcorr}[k] = \mathrm{xcorr}_x[k] + \mathrm{xcorr}_y[k]$,
or equivalently by using the complex-valued cardinal embedding
$\Cardinal_\mathbb{C}(s) = x + iy$ directly in equation~\eqref{eq:xcorr}.
\end{remark}

\begin{definition}[Normalized Cross-Correlation Score]
\label{def:normalized_xcorr}
The \emph{normalized cross-correlation} between $q$ and $t$ at best alignment
offset $k^*$ is:
\begin{equation}
  \rho(q, t) = \frac{\mathrm{xcorr}[k^*]}{\sqrt{\sum_i |\mathbf{q}_i|^2 \cdot \sum_j |\mathbf{t}_j|^2}},
  \qquad k^* = \arg\max_k |\mathrm{xcorr}[k]|,
  \label{eq:normalized_xcorr}
\end{equation}
where the denominator normalizes by the geometric mean of sequence energies.
\end{definition}

\begin{theorem}[Matched Filter Optimality]
\label{thm:matched_filter}
The matched filter cross-correlation of Definition~\ref{def:cross_correlation}
is the optimal (minimum mean-square error, maximum signal-to-noise ratio)
detector for query sequence $q$ embedded at an unknown offset in target $t$,
under the model in which the target outside the match is modeled as
white noise in cardinal space.
\end{theorem}

\begin{proof}
This is the Wiener--Hopf result from linear detection theory, applied to the
cardinal-embedded sequences.

\medskip
\noindent\textbf{Model.} Let $t[k] = q[k - k^*] + n[k]$ where $q[k-k^*]$ is
the query embedded at offset $k^*$ and $n[k]$ is additive noise with power
spectral density $S_n(\omega)$. We seek the linear filter $h$ such that the
filtered output $y = h * t$ maximizes the signal-to-noise ratio (SNR) at $k=k^*$.

\medskip
\noindent\textbf{Frequency-domain SNR.} The SNR at the output of a filter $H(\omega)$
applied to the noisy input is:
\begin{equation}
  \mathrm{SNR} = \frac{|H(\omega) Q(\omega)|^2}{|H(\omega)|^2 S_n(\omega)}
  = \frac{|Q(\omega)|^2}{S_n(\omega)},
\end{equation}
where $Q(\omega) = \mathrm{FFT}(\mathbf{q})(\omega)$.
By the Cauchy--Schwarz inequality:
\begin{equation}
  \left|\int H(\omega) Q(\omega)\,d\omega\right|^2
  \leq \int |H(\omega)|^2 S_n(\omega)\,d\omega \cdot \int \frac{|Q(\omega)|^2}{S_n(\omega)}\,d\omega,
\end{equation}
with equality when $H(\omega) = c \, \overline{Q(\omega)} / S_n(\omega)$.
For white noise, $S_n(\omega) = \sigma^2$ (constant), and the optimal filter
reduces to $H(\omega) = c \, \overline{Q(\omega)}$, which in the time domain
is the cross-correlation with the time-reversed (conjugated) query.

\medskip
\noindent\textbf{Implementation.} The matched filter output at lag $k$ is:
\begin{equation}
  y[k] = (h * t)[k] = \sum_j \overline{q[j]} \cdot t[k+j]
  = \mathrm{xcorr}(q,t)[k].
\end{equation}
This is exactly the cross-correlation of equation~\eqref{eq:xcorr},
confirming that the matched filter is implemented by cross-correlation.
FFT-based computation achieves $O(T \log T)$ complexity rather than the
naive $O(LT)$ of sliding-window convolution. \qed
\end{proof}

\begin{remark}[Physical Interpretation]
In optical coherence: two coherent waves interfere constructively when their
phase difference is zero (path lengths differ by an integer number of wavelengths)
and destructively when the phase difference is $\pi$. The cross-correlation
peak at offset $k^*$ corresponds to constructive interference of the query
waveform with the target at alignment position $k^*$. All other offsets
produce near-zero cross-correlation (destructive interference of a periodic
waveform with a shifted copy of itself). This is the same physics that produces
Young's double-slit fringes; the nucleotide sequence is the wave.
\end{remark}

%------------------------------------------------------------------------------
\section{Phase-Preserving vs.\ Magnitude-Only Analysis}
\label{sec:phase_vs_magnitude}
%------------------------------------------------------------------------------

Chapter~\ref{chap:shader_spectrometer} computed only magnitude spectra
$|X_k|, |Y_k|$, discarding phase information. The matched filter retains both
magnitude and phase. The difference in peak width is quantitatively explained
by the Nyquist sampling theorem.

\begin{theorem}[Phase-Preserving Localization]
\label{thm:phase_localization}
Let $q$ and $t$ be identical sequences of length $L$ (query is an exact copy
of a segment of the target). Then:
\begin{enumerate}[label=(\roman*)]
  \item The phase-preserving matched filter cross-correlation
        $|\mathrm{xcorr}(q,t)[k]|$ achieves a peak width of at most
        $2$ base pairs (half-maximum full-width at the 3\,dB point).
  \item The magnitude-only cross-correlation $\mathrm{xcorr}(|\Phi_K(q)|, |\Phi_K(t)|)$
        achieves a peak width of approximately $L/K_{\max}$ base pairs,
        which for $K_{\max} = 12$ and $L = 200$ is $\approx 17$\,bp minimum
        (observed 100--130\,bp due to aliasing from mode truncation).
\end{enumerate}
\end{theorem}

\begin{proof}
\textbf{(i).} For $t = q$ (exact match at offset $k^* = 0$):
\begin{equation}
  \mathrm{xcorr}(q,q)[k] = \mathrm{IFFT}\bigl(|\mathrm{FFT}(\mathbf{q})|^2\bigr)[k],
\end{equation}
which is the autocorrelation function of $\mathbf{q}$. For a sequence of length
$L$ with unit-variance coefficients, the autocorrelation is $L$ at $k=0$ and
$O(\sqrt{L})$ at $|k| \geq 1$ by the central limit theorem applied to the sum
$\sum_i q_i q_{i+k}$ for $k \neq 0$. The half-maximum point (3\,dB down) occurs
at $|k| = 1$, since $|\mathrm{xcorr}[1]|/|\mathrm{xcorr}[0]| \to 0$ for
large $L$ and random-phase sequence. The peak width is therefore $2$\,bp
(from $k=-1$ to $k=+1$ at the 3\,dB level), which equals $1/f_{\mathrm{Nyquist}}$
in base-pair units. This is the Nyquist limit: the autocorrelation of a
length-$L$ discrete sequence has the minimum resolvable peak width of 1 sample
= $1$\,bp, and the observed 3\,dB width of 2\,bp (2 samples) is the standard
result for a bandlimited signal.

\textbf{(ii).} In magnitude-only analysis, phase information is destroyed.
The cross-correlation of the magnitude spectra $\Phi_K(q) = (|X_1|,\ldots,|X_K|)$
is a $K$-dimensional dot product rather than an $L$-dimensional cross-correlation.
The effective bandwidth of the representation is limited to $K=12$ modes,
corresponding to wavelengths $\geq L/K = 200/12 \approx 17$\,bp.
By the bandwidth--resolution duality (time-bandwidth product), a bandwidth
of $K/L$ in normalized frequency corresponds to a temporal resolution of
$L/K \approx 17$\,bp. In practice, the observed peak width of 100--130\,bp
exceeds this theoretical minimum because (a) mode truncation at $K=12$ removes
high-frequency information needed to sharpen the peak, and (b) aliasing from
the abrupt bandwidth cutoff creates sidelobes that raise the apparent baseline.
The factor of $\approx 7$--$8$ broadening beyond the theoretical minimum is
consistent with the Gibbs phenomenon for a rectangular spectral window. \qed
\end{proof}

\begin{keyresult}[Peak Width Summary]
\begin{center}
\begin{tabular}{lcc}
\toprule
Method & Phase & Peak Width (bp) \\
\midrule
Phase-preserving matched filter & Full & $\approx 2$ \\
Magnitude-only (Ch.\ 26) & Discarded & 100--130 \\
Theoretical minimum (magnitude) & Discarded & $\approx 17$ \\
\bottomrule
\end{tabular}
\end{center}
The 2\,bp peak width equals the Nyquist limit; the 100--130\,bp broadening
is explained by the Gibbs phenomenon from rectangular spectral windowing.
\end{keyresult}

%------------------------------------------------------------------------------
\section{Statistical Calibration via z-Score}
\label{sec:z_score}
%------------------------------------------------------------------------------

The matched filter operates on a target of length $T$ and produces $T$
cross-correlation values. A detection threshold is needed. The $z$-score
provides optimal calibration when the background distribution is normal.

\begin{proposition}[Null Distribution Normality]
\label{prop:null_normal}
Under the null hypothesis (query $q$ does not appear in target $t$, and the
target nucleotide sequence is approximately i.i.d.\ uniform), the cross-correlation
values $\{\mathrm{xcorr}[k] : k \neq k^*\}$ are approximately normally distributed
with mean $\mu_0 = 0$ and standard deviation $\sigma_0 = O(\sqrt{L})$.
\end{proposition}

\begin{proof}
For $k \neq k^*$, the cross-correlation is $\mathrm{xcorr}[k] = \sum_{j=1}^L q_j t_{k+j}$.
Under the null, $t_{k+j}$ is i.i.d.\ with $\mathbb{E}[t_{k+j}] = 0$ (balanced
nucleotide composition gives zero mean in cardinal coordinates) and
$\mathrm{Var}(t_{k+j}) = 1/2$ (each of the four cardinal values has equal
probability, and the mean square of the x- and y-coordinates is $1/2$).
By the central limit theorem, $\mathrm{xcorr}[k] = \sum_j q_j t_{k+j}$ is
approximately $\mathcal{N}(0, L/2)$ for large $L$ (since $q_j$ are fixed
constants with $\sum_j q_j^2 = L/2$ and $t_{k+j}$ are i.i.d.). Thus
$\sigma_0 = \sqrt{L/2}$, and the null distribution is normal. \qed
\end{proof}

\begin{definition}[z-Score Detection Statistic]
\label{def:zscore}
The \emph{$z$-score} of the peak cross-correlation value is:
\begin{equation}
  z = \frac{\max_k |\mathrm{xcorr}[k]| - \mu_0}{\sigma_0},
  \label{eq:zscore}
\end{equation}
where $\mu_0$ and $\sigma_0$ are the mean and standard deviation of the
cross-correlation values computed empirically from the same target sequence.
\end{definition}

\begin{theorem}[F1 Optimality at z $\geq$ 6]
\label{thm:f1_optimal}
The detection threshold $z_{\mathrm{thresh}} = 6$ achieves F1 $= 1.000$
in the validated setting of \citet{Sachikonye2025_Interference}.
\end{theorem}

\begin{proof}
By Proposition~\ref{prop:null_normal}, the null distribution is
$\mathcal{N}(\mu_0, \sigma_0^2)$. The probability of a false positive at
threshold $z=6$ is $P(|Z| \geq 6) \approx 2 \times 10^{-9}$ for a standard
normal $Z$. For a 1\,Mb target with $\approx 10^6$ windows, the expected number
of false positives is $\approx 2 \times 10^{-3}$---less than one in a thousand
queries. Simultaneously, the true positive signal at the planted motif position
achieves $z \gg 6$ (empirically $z > 20$ at 0\% mutation; $z > 8$ at 50\%
mutation as measured in \citet{Sachikonye2025_Interference}), so the true
positive is detected with probability 1. Therefore precision $= 1.000$ and
recall $= 1.000$, giving F1 $= 1.000$. \qed
\end{proof}

%------------------------------------------------------------------------------
\section{Multi-Channel Analysis and Protein Extension}
\label{sec:multichannel}
%------------------------------------------------------------------------------

The two-channel cardinal embedding for DNA naturally extends to amino acid
sequences by choosing biochemically meaningful channels.

\begin{definition}[Protein Cardinal Embedding]
\label{def:protein_embedding}
For amino acid sequences, the two-channel embedding uses:
\begin{enumerate}[label=(\roman*)]
  \item \textbf{Channel 1}: Kyte--Doolittle hydropathy index $h(a_i)$,
        normalized to $[-1, +1]$.
  \item \textbf{Channel 2}: Van der Waals volume $v(a_i)$,
        normalized to $[-1, +1]$.
\end{enumerate}
The matched filter cross-correlation is then applied to $(h(a_1),\ldots,h(a_L))$
and $(v(a_1),\ldots,v(a_L))$ independently, with results combined by addition.
\end{definition}

\begin{remark}[Biochemical Interpretation]
Hydropathy encodes the burial tendency of each residue (membrane topology);
van der Waals volume encodes steric packing. Together they capture the two
most structurally determining properties of an amino acid sequence. The
matched filter on these channels detects structural (fold-level) similarity
rather than raw sequence identity, implementing the partition-theoretic principle
that functional identity = partition-topological identity
(Chapter~\ref{chap:enzymes_topology}).
\end{remark}

%------------------------------------------------------------------------------
\section{The Planted Motif Validation}
\label{sec:planted_motif}
%------------------------------------------------------------------------------

The matched filter interference method was validated on synthetic 1\,Mb loci
with controlled planted features \citep{Sachikonye2025_Interference}.

\begin{keyresult}[Planted Motif Recovery]
A synthetic 1\,Mb locus was constructed with:
\begin{itemize}
  \item A canonical 200\,bp motif planted at position 500,000\,bp.
  \item A 10\%-mutated copy of the same motif at position 800,000\,bp.
  \item Background sequence with 50\% additional random mutation.
\end{itemize}
Results:
\begin{itemize}
  \item Canonical motif (position 500,000): recovered at rank 1 with $\rho = 0.9969$.
  \item Mutated copy (position 800,000): recovered at rank 2 with $\rho = 0.8979$.
  \item Both recovered with F1 $= 1.000$ at $z$-score threshold $\geq 6$.
  \item Recovery held across all 60 tested positions (60/60 at 50\% mutation).
\end{itemize}
\end{keyresult}

These results are the quantitative statement that sequence alignment is wave
interference detection: constructive interference at the correct alignment
offset produces a peak ($\rho = 0.9969$) while background mutation produces
near-zero correlation at all other positions.

%------------------------------------------------------------------------------
\section{Validated Performance at All Mutation Rates}
\label{sec:interference_performance}
%------------------------------------------------------------------------------

\begin{keyresult}[AUC and Position Recovery]
Results from \citet{Sachikonye2025_Interference}:
\begin{center}
\begin{tabular}{ccc}
\toprule
Mutation Rate & AUC & Positions Recovered (of 60) \\
\midrule
0\%  & 1.000 & 60/60 \\
10\% & 1.000 & 60/60 \\
20\% & 1.000 & 60/60 \\
30\% & 1.000 & 60/60 \\
40\% & 1.000 & 60/60 \\
50\% & 1.000 & 60/60 \\
\bottomrule
\end{tabular}
\end{center}
AUC $= 1.000$ is maintained at all tested mutation rates. The matched filter
degrades gracefully: at 50\% mutation, the peak is broader and lower, but still
detectable above $z=6$ because the background cross-correlation remains
normally distributed with the same $\sigma_0$.
\end{keyresult}

\begin{keyresult}[Computational Performance]
Measured in a JavaScript implementation \citep{Sachikonye2025_Interference}:
\begin{itemize}
  \item Preprocessing (precompute FFT of 1\,Mb target): $\approx 530$\,ms.
  \item Per-query scan (1\,Mb target): $\approx 836$\,ms.
  \item Complexity: $O(T \log T)$ for FFT preprocessing; $O(T)$ per query
        (IFFT of the pointwise product of precomputed target FFT and query FFT).
  \item Scalability: preprocessing amortized over all queries to the same target.
\end{itemize}
\end{keyresult}

%------------------------------------------------------------------------------
\section{Constructive and Destructive Interference: The Partition Interpretation}
\label{sec:partition_interference}
%------------------------------------------------------------------------------

The wave interference interpretation has a precise partition-theoretic reading.

\begin{proposition}[Constructive Interference = Partition Boundary Alignment]
\label{prop:constructive_interference}
The cross-correlation peak at offset $k^*$ corresponds to the alignment of
partition boundaries in the cardinal waveforms of $q$ and $t$.
\end{proposition}

\begin{proof}
In the cardinal representation, a nucleotide is a charge assignment: $+1$ or
$-1$ in one channel, $0$ in the other (see equation~\eqref{eq:cardinal_map}).
A transition between nucleotides (e.g.\ $\mathtt{A} \to \mathtt{G}$) corresponds
to a sign change or a non-zero-to-zero transition in the waveform, i.e., a
partition boundary (Chapter~\ref{chap:charge_emergence}: partition boundaries
carry charge). The cross-correlation at lag $k$:
\begin{equation}
  \mathrm{xcorr}[k] = \sum_{j=1}^L q_j t_{k+j}
\end{equation}
is large (constructive interference) when $q_j$ and $t_{k+j}$ have the same
sign pattern---meaning their partition boundaries align. Misaligned boundary
patterns give cancellation (destructive interference) since positive and
negative contributions to the sum are balanced. The offset $k^*$ at which
partition boundaries align is exactly the correct sequence alignment, confirming
that alignment detection is partition boundary detection. \qed
\end{proof}

\begin{corollary}[Sequence Alignment = Partition Boundary Detection]
\label{cor:alignment_boundary}
Optimal sequence alignment is the problem of finding the offset at which
the cardinal partition boundaries of query and target constructively interfere.
The matched filter is the optimal detector for this offset.
\end{corollary}

This corollary unifies sequence alignment (a computational problem) with
partition boundary detection (a physical measurement). They are the same
operation in different coordinate systems.

%------------------------------------------------------------------------------
\section{Connection to Optical Interference}
\label{sec:optical_connection}
%------------------------------------------------------------------------------

The formal analogy between sequence cross-correlation and optical interference
is exact at the level of the mathematics.

\begin{proposition}[Formal Identity with Optical Cross-Correlation]
\label{prop:optical_identity}
Let $E_1(t)$ and $E_2(t)$ be two coherent optical fields. The optical intensity
at the beam-splitter output is:
\begin{equation}
  I(\tau) = |E_1(t) + E_2(t+\tau)|^2 = I_1 + I_2 + 2\,\mathrm{Re}[\Gamma_{12}(\tau)],
\end{equation}
where $\Gamma_{12}(\tau) = \langle E_1^*(t) E_2(t+\tau)\rangle$ is the
mutual coherence function. For ergodic fields, $\Gamma_{12}(\tau) = \mathrm{xcorr}(E_1,E_2)[\tau]$.
This is equation~\eqref{eq:xcorr} with $E_1 = \mathbf{q}$ and $E_2 = \mathbf{t}$.
\end{proposition}

\begin{proof}
Expanding $|E_1 + E_2|^2 = |E_1|^2 + |E_2|^2 + 2\,\mathrm{Re}[E_1^* E_2]$
and taking the expectation gives the van Cittert--Zernike expression. For
ergodic processes (time average = ensemble average), the cross-term
$\langle E_1^*(t) E_2(t+\tau)\rangle$ equals the cross-correlation function.
Applying to cardinal-embedded sequences: $E_1 \leftrightarrow \mathbf{q}$,
$E_2 \leftrightarrow \mathbf{t}$, $\tau \leftrightarrow k$ (discrete lag).
The cross-correlation of equation~\eqref{eq:xcorr} is the mutual coherence
function of the two nucleotide waveforms. \qed
\end{proof}

The van Cittert--Zernike theorem from optical coherence theory is, in the
partition framework, the statement that the spectral alignment of two cardinal
waveforms measures their partition boundary coherence. Constructive interference
(large $|\Gamma_{12}(k^*)|$) signals that the two sequences share the same
partition boundary structure at offset $k^*$---they are the same sequence,
possibly mutated.

%------------------------------------------------------------------------------
\section{Scaling and Implementation}
\label{sec:interference_scaling}
%------------------------------------------------------------------------------

\begin{proposition}[Linear Scaling with Target Length]
\label{prop:linear_scaling}
The matched filter analysis of a query of length $L$ against a target of
length $T \geq L$ scales as:
\begin{enumerate}[label=(\roman*)]
  \item \textbf{Preprocessing} (precompute target FFT): $O(T \log T)$.
  \item \textbf{Per-query scan}: $O(L \log L)$ (query FFT) $+$ $O(T)$ (pointwise product and IFFT).
  \item \textbf{Total for $Q$ queries}: $O(T \log T + Q(L \log L + T))$.
\end{enumerate}
\end{proposition}

\begin{proof}
The FFT of a length-$T$ sequence requires $O(T \log T)$ operations. Once the
target FFT $\hat{\mathbf{t}} = \mathrm{FFT}(\mathbf{t})$ is precomputed, the
cross-correlation for a new query is computed as:
(1) FFT of query: $O(L \log L)$; (2) pointwise product
$\overline{\hat{\mathbf{q}}} \cdot \hat{\mathbf{t}}$: $O(T)$; (3) IFFT: $O(T \log T)$.
Steps (2) and (3) dominate, giving $O(T)$ per query after preprocessing (or
$O(T \log T)$ if the IFFT must be computed at full length $T$). The total
scales linearly with $T$, confirming the implementation measurements. \qed
\end{proof}

\begin{remark}[Comparison with Smith-Waterman]
Smith-Waterman dynamic programming scales as $O(LT)$ for a query of length $L$
against a target of length $T$. For $L = 200$, $T = 10^6$: Smith-Waterman
requires $2 \times 10^8$ operations; the matched filter requires $O(T \log T)
\approx 2 \times 10^7$ operations (a $10\times$ speedup from the FFT algorithm
alone). Combined with the precomputation amortized over multiple queries and the
superior sensitivity at high mutation rates (AUC $= 1.000$ vs. Smith-Waterman
AUC degrading at $> 30\%$ mutation), the matched filter is the preferred
implementation of the optimal sequence alignment method.
\end{remark}

%------------------------------------------------------------------------------
\section{The Interference Pattern as a Partition Map}
\label{sec:interference_partition_map}
%------------------------------------------------------------------------------

The full cross-correlation vector $\mathrm{xcorr}(q,t)[k]$ for $k = 0,\ldots,T-1$
is not merely a search output. It is a map of the partition structure of the target.

\begin{definition}[Partition Map of Target]
\label{def:partition_map}
For a fixed query $q$ and target $t$, the \emph{partition map} is the function:
\begin{equation}
  \mathcal{M}_{q,t} : \{0, 1, \ldots, T-1\} \to \mathbb{R},
  \qquad k \mapsto |\mathrm{xcorr}(q,t)[k]|.
  \label{eq:partition_map}
\end{equation}
High values of $\mathcal{M}_{q,t}(k)$ indicate positions where the target
contains a partition boundary structure similar to that of the query.
\end{definition}

\begin{proposition}[Partition Map Detects All Homologous Loci]
\label{prop:partition_map_detects}
For a target containing $M$ copies of $q$ (exact or mutated) at positions
$k_1^*, \ldots, k_M^*$, the partition map $\mathcal{M}_{q,t}$ has local maxima
at $k_1^*, \ldots, k_M^*$ with values $\rho_m = \rho(q, t_{k_m^*})$ where
$\rho$ is the normalized cross-correlation of equation~\eqref{eq:normalized_xcorr}.
\end{proposition}

\begin{proof}
At each position $k_m^*$, the target segment $t_{k_m^*:k_m^*+L}$ is a mutated
copy of $q$. The cross-correlation at this lag is the inner product of $q$ with
a similar waveform, producing a large value. At positions $k \neq k_m^*$
(background), the target segment is uncorrelated with $q$, giving
$|\mathrm{xcorr}[k]| = O(\sqrt{L})$ (background level) rather than $O(L)$
(signal level). Therefore $\mathcal{M}_{q,t}$ has local maxima at the planted
positions, with heights proportional to the similarity of the motif copy to
the query. \qed
\end{proof}

The partition map $\mathcal{M}_{q,t}$ is the genomic analogue of the
interference pattern in a Michelson interferometer: it maps all positions where
the query waveform is coherent with the target, providing a comprehensive view
of all homologous regions.

%------------------------------------------------------------------------------
\section*{Chapter Summary}
%------------------------------------------------------------------------------

\begin{chapterbox}[title=Chapter 27 --- Key Results Established]
\begin{enumerate}[label=\textbf{\arabic*.}]
  \item \textbf{Sequence Matching = Wave Interference}
        (Proposition~\ref{prop:matching_is_interference}): Under the cardinal
        embedding, sequence alignment is signal detection of a known waveform
        in a noisy background.

  \item \textbf{Matched Filter Optimality}
        (Theorem~\ref{thm:matched_filter}): The cross-correlation
        $\mathrm{IFFT}(\overline{\mathrm{FFT}(q)} \cdot \mathrm{FFT}(t))$ is the
        Wiener--Hopf optimal detector for cardinal-embedded sequences.

  \item \textbf{Phase-Preserving Localization}
        (Theorem~\ref{thm:phase_localization}): Phase-preserving matched filter
        achieves 2\,bp peak width (Nyquist limit); magnitude-only achieves
        100--130\,bp (Gibbs broadening at $K=12$).

  \item \textbf{Null Distribution Normality} (Proposition~\ref{prop:null_normal}):
        Background cross-correlation is $\mathcal{N}(0, L/2)$, validating
        $z$-score calibration.

  \item \textbf{F1 Optimality at z $\geq$ 6} (Theorem~\ref{thm:f1_optimal}):
        False-positive probability $\approx 2\times10^{-9}$; F1 $= 1.000$.

  \item \textbf{Validated Performance} \citep{Sachikonye2025_Interference}:
        60/60 position recovery and AUC $= 1.000$ at all mutation rates
        0\%--50\%; planted motif $\rho = 0.9969$ (rank 1) and $\rho = 0.8979$
        (rank 2) recovered through 50\% background mutation.

  \item \textbf{Constructive Interference = Partition Boundary Alignment}
        (Proposition~\ref{prop:constructive_interference}): The cross-correlation
        peak signals alignment of partition boundaries, identifying sequence
        alignment with partition boundary detection.

  \item \textbf{Formal Identity with Optics}
        (Proposition~\ref{prop:optical_identity}): The cardinal cross-correlation
        is the van Cittert--Zernike mutual coherence function, completing the
        identification of sequence matching with optical interference.

  \item \textbf{Computational Scaling} (Proposition~\ref{prop:linear_scaling}):
        $O(T \log T)$ preprocessing, $O(T)$ per query; measured at 530\,ms
        (preprocess 1\,Mb) and 836\,ms (query 1\,Mb).
\end{enumerate}
\end{chapterbox}
